\subsection{The geometric mean is a Haar mean}
\label{sec:curvature-haar-mean}

We now identify the mean defined by the original Hodge metric, once
its complex uniformization and local metric identifications are given.
For the matrix domain
\[
 \mathcal D_{p,q}=\{Z\in M_{p,q}(\mathbf C):I_q-Z^*Z>0\},
 \qquad \Phi(Z)=-\log\det(I_q-Z^*Z),
\]
the determinant quotient Hodge norm has local logarithm $\Phi/2$,
up to a pluriharmonic term. Indeed, on the positive plane
$\operatorname{graph}(W)$ the quotient metric has Gram matrix
$(I_q-WW^*)^{-1}$; its determinant norm is the square root of
its determinant. The transpose coordinate $Z=W^t$ gives the
displayed potential.

Put $\Phi=\sum_\sigma\Phi_\sigma$ on the product domain and
$n=\sum_\sigma p_\sigma q_\sigma$. Direct differentiation gives
\[
 (\partial\bar\partial\Phi_\sigma)_Z(T,\bar T)
 =\operatorname{tr}\bigl((I-ZZ^*)^{-1}
                       T(I-Z^*Z)^{-1}T^*\bigr).
\]
The determinant of this Hessian is
$\det(I-Z^*Z)^{-(p_\sigma+q_\sigma)}$. Therefore, with
$dd^c=(i/\pi)\partial\bar\partial$,
\[
 (dd^c(\Phi/2))^n=c_n b\,dV,\qquad
 c_n=\frac{n!}{\pi^n},\qquad
 b=\prod_\sigma\det(I-Z_\sigma^*Z_\sigma)^{-(p_\sigma+q_\sigma)}.
\]
This includes zero-dimensional factors by the empty-determinant
convention. In our equal-signature case, $b=\exp(N\Phi)$.

The fractional linear unitary action changes $\Phi$ by
$2\log|\det(CZ+D)|$. Its Levi form is unchanged, so $b\,dV$
is invariant. Since the stabilizer $K$ is compact, the pushforward
of Haar measure under $G\to G/K$ is $c\,b\,dV$ for some
$0<c<\infty$.

\begin{proposition}\label{prop:geometric-haar-mean}
Retain the original finite collection of arithmetic components of
complex dimension $n$, their measurable fundamental sets $\mathcal F_k$,
and restrictions $\gamma_m:B_m\to X$ of local holomorphic sheets, with
$B_m\subset\mathcal F_{k(m)}$. Suppose these sheets are disjoint
and cover up to chart-null sets, their multiplicity over almost
every point of $\mathcal F_k$ is $d_k\in\mathbf Z_{\ge0}$ with
$\sum_k d_k>0$, and their original Hodge
potential is $\Phi_k/2+\Re H_m$ locally. Then the original
normalized geometric mean is
\[
 \operatorname{mean}_X(g)=
 \sum_{k:d_k>0}\frac{d_kV_k}{\sum_\ell d_\ell V_\ell}
                 \int\bar g_k\,d\nu_k,\qquad
 V_k=\int_{\mathcal F_k}b_k\,dV,
\]
where $\nu_k$ is normalized arithmetic Haar measure and
$\bar g_k$ averages the $d_k$ sheet values. The assertion holds
for every absolutely integrable function $g$ on the projective variety.
\end{proposition}

\begin{proof}
Holomorphic change of variables gives the original analytic
integral as
\[
 I(g)=c_n\sum_m\int_{B_m}g(\gamma_m(t))b_{k(m)}(t)\,dV(t).
\]
Use nonnegative measures first; absolute integrability then
justifies signed sums. The arithmetic finite-volume theorem gives
$0<V_k<\infty$. Consequently
$I(1)=c_n\sum_k d_kV_k>0$, and division gives the formula.
The original denominator is defined by this same Hodge-potential
integral against $1$, so it is exactly $I(1)$.

For precision, let $s_k$ be the order of the kernel of
$\Gamma_k$ on its domain, and write
$(\pi_k)_*dg_k=c_kb_k\,dV$. Almost every point has stabilizer
of order $s_k$: every additional fixed locus is a proper analytic
subset. The raw quotient volume is
$(c_k/s_k)V_k$. Thus normalized component Haar measure is
$b_k\,dV/V_k$, while the weights retain $d_kV_k$.
No stabilizer or relative volume is suppressed.
\end{proof}

For lattice invariants the prescribed fibre isometries make all
sheet values equal. Proposition~\ref{prop:arithmetic-isotropic-integrability}
therefore gives finite mean and the same radius scaling for the
original isotropic count. The cusp bound applies with its original
normalization. The required geometric input is the actual
uniformization with these sheets and this metric; their existence
is not inferred from a probability-measure contract. The arithmetic
model and signed boundary intersection limit remain separate.
